% Første linje er dokumentklassen. Brug *altid* memoir!
\documentclass[a4paper,oneside,article]{memoir}

% Herunder indlæses forskellige pakker
\usepackage[utf8]{inputenc}  % Korrekt håndtering af æ, ø og å
\usepackage[T1]{fontenc}  % Korrekt håndtering af æ, ø og å
\usepackage{microtype}  % Typografisk magi! Giver bl.a. pænere orddeling
\usepackage{graphicx}  % Gør det muligt at indsætte billeder
\usepackage{amsmath}  % Giver adgang til uundværlige matematikting
\usepackage{mathtools}  % Retter ting i ams og giver ekstra muligheder
\usepackage[danish]{babel}  % Danske betegnelser og orddeling
\usepackage{float} %Mere kontrol over billede-placering
\renewcommand{\danishhyphenmins}{22}  % Bedre dansk orddeling
\usepackage[exponent-product=\ensuremath{\cdot}]{siunitx} % Sientific notation for forskellige ting, med en regel om at den skal bruge cdot istedet for times. 
\usepackage{derivative} % Til at skrive differentialligninger 
\usepackage{xcolor} % For at få flere tekst farver 
\usepackage[colorlinks=true,pdftex]{hyperref} % For at få refrences man kan trykke på, slet [] hvis man vil have kasser i stedet for farvet tekst.
\usepackage{pdfpages} % Til at indsætte andre pdf'er
\usepackage{lmodern} % Så man kan compile i TeXworks
\usepackage{alphabeta} % Til at skrive græske bogstaver
\usepackage{subfig} % For at have to billeder ved siden af hinanden

% Pænere toc indents 
\usepackage{tocloft}
\cftsetindents{section}{1em}{2em}
\cftsetindents{subsection}{4em}{0em}


% Indholdet i dit dokument starter her!
\begin{document}
\author{Mikkel Moth Billing \\ 202307989}
\title{Title}
\date{\today}
%\maketitle   % Indsætter overskriften
% Table of Contents indstillinger
\setcounter{tocdepth}{2}
\setcounter{secnumdepth}{1}
\hypersetup{linkcolor=black} % Laver tableofcontents sort.  
%\tableofcontents %Indsætter ToC
\hypersetup{linkcolor=red} % og så alle andre hyperrefs røde. 

%Custom Commands
%Til opgavebeskrivelser
\newcommand{\opg}[1]{\textcolor{gray}{\emph{#1}}\plainbreak{1}} 

\newcommand{\ti}[1]{\cdot 10^{#1}} % Let: gange 10^x
\newcommand{\e}[1]{\emph{e}^{#1}} % Pænere e^x

% Lette paranteser 
\newcommand{\br}[2][1]{%
    \ifcase#1
    \or \left( #2 \right) %1 eller ingenting
    \or \left[ #2 \right] %2
    \or \left\langle #2 \right\rangle %3
    \or \left\{ #2 \right\} %4
    \or \left| #2 \right| %5
    \fi
}
%Til at sætter bileder ind i store dokumenter
\newcommand{\bil}[1]{
\begin{figure}[H]
    \centering
    \includegraphics[width=0.75\linewidth]{Billeder/#1.png}
    \label{fig:#1}
\end{figure}
}

%Til at sætte pdf'er ind i store dokumenter (kun pdf'er på mere end 1 side)
\newcommand{\pdf}[3]{
\includepdf[pages=#2,scale=1,pagecommand={\subsection{#1}  \label{asc:#2}}]{pdf/emo.pdf}
\includepdf[pages=\the\numexpr 1 + #2\relax-#3, scale=1]{pdf/emo.pdf}
}

\newcommand{\side}[2]{
\includepdf[pages=#2,scale=.85,pagecommand={\subsection{#1}  \label{asc:#1}}]{pdf/emo.pdf}
}

% Lav ny pagestyle (kun odd fordi enkeltsidet)
\makepagestyle{Title}
\pagestyle{Title}
\makeoddhead{Title}{\footnotesize\color{gray}{\theauthor}}{}{\textcolor{gray}{\thedate}}
\makeoddfoot{Title}{}{\footnotesize\color{gray}{\thepage{} / \thelastpage}}{}
% Også fod på siden med \maketitle
\makeoddfoot{plain}{}{\footnotesize\color{gray}{\thepage{} / \thelastpage}}{}
\newpage

%Skriv
\chapter{Elektricitet}
\side{Elektrisk ladning og felter}{736}

\side{Gauss's lov}{768}

\side{Elektrisk potentiale}{800}




\end{document} % ----------Dokumentet er nu slut! ------------